%!TeX program = xelatex
\documentclass{SYSUReport}

% 根据个人情况修改
\headl{实验零：实验准备}
\headr{操作系统原理实验}
\lessonTitle{操作系统原理实验}
\reportTitle{实验零：实验准备}
\inst{计算机学院}
\major{保密管理}
\stuname{吕梓文}
\stuid{24353028}
\date{\today}

\AddEnumerateCounter{\chinese}{\chinese}{}
\begin{document}

\cover
\updateheader{实验零：实验准备}
\thispagestyle{empty} % 首页不显示页码
\clearpage

% \pagenumbering{Roman}
% \setcounter{page}{1}
% \tableofcontents
\clearpage

\pagenumbering{arabic}
\setcounter{page}{1}

%%可选择这里也放一个标题
%\begin{center}
%    \title{ \Huge \textbf{{标题}}}
%\end{center}

\section{实验零}
\updateheader{实验零：实验准备}
\begin{enumerate}[label={\chinese*、},labelsep=0pt]
    
    \item \textbf{Linux 下环境配置}
    
由于我有一台 Linux 实体机，系统为 Ubuntu 24.04 LTS，镜像源已经配置好了，所以这里就不再赘述

    \begin{enumerate}[label = \arabic*、]
        \item 安装依赖
        
        \begin{minted}[linenos=false]{bash}
        sudo apt update && sudo apt upgrade
        sudo apt install -y \
            build-essential \
            gdb \
            qemu-system-x86
        \end{minted}

        \item 为 \verb|qemu-system-x86\_64| 这个包名字太长，设置一个别名
        
        \begin{minted}[linenos=false]{bash}
        echo "alias qemu='qemu-system-x86_64'"    
        \end{minted}


    \end{enumerate}
    
    \item \textbf{Rust 编程练习}
    
    \textbf{任务1：基础文件操作与倒计时}
    
    实现了三个函数：
    \begin{itemize}
        \item \texttt{count\_down(seconds: u64)}：倒计时函数，每秒输出剩余秒数
        \item \texttt{read\_and\_print(file\_path: \&str)}：读取并打印文件内容
        \item \texttt{file\_size(file\_path: \&str) -> Result<u64, \&str>}：获取文件大小
    \end{itemize}
    
    核心代码实现（\texttt{utils.rs}）：
    \begin{minted}{rust}
    pub fn count_down(seconds: u64) {
        for i in (1..=seconds).rev() {
            thread::sleep(Duration::from_secs(1));
            println!("{i} seconds left");
        }
        println!("Countdown finished!");
    }

    pub fn read_and_print(file_path: &str) -> io::Result<()> {
        let data = fs::read_to_string(file_path)?;
        println!("{}", data);
        Ok(())
    }

    pub fn file_size(file_path: &str) -> Result<u64, &'static str> {
        let metadata = fs::metadata(file_path)
            .map_err(|_| "File not found!")?;
        Ok(metadata.len())
    }
    \end{minted}
    
    主函数调用逻辑（\texttt{main.rs}）：
    \begin{minted}{rust}
    mod utils;
    use std::io;
    use utils::{count_down, file_size, read_and_print};

    fn main() {
        count_down(5);
        if let Err(_) = read_and_print("/etc/hosts") {
            println!("File not found!");
        }

        loop {
            let mut input = String::new();
            if let Err(e) = io::stdin().read_line(&mut input) {
                println!("{}", e.to_string());
            }
            let path = input.trim();
            if path == "q" { break; }
            match file_size(path) {
                Err(e) => println!("{}", e),
                Ok(sz) => println!("The size of the file {} is {}", path, sz),
            }
        }
    }
    \end{minted}
    
    \textbf{任务2：字节数转换函数}
    
    实现了 \texttt{humanized\_size} 函数，将字节数转换为人类可读格式：
    \begin{minted}{rust}
    pub fn humanized_size(size: u64) -> (f64, &'static str) {
        let units = vec!["B", "KiB", "MiB", "GiB", "TiB"];
        let mut float_size = size as f64;
        let mut cnt = 0;
        
        while float_size > 1024.0 && cnt < units.len() - 1 {
            float_size /= 1024.0;
            cnt += 1;
        }
        (float_size, units[cnt])
    }

    #[test]
    fn test_humanized_size() {
        let byte_size = 1554056;
        let (size, unit) = humanized_size(byte_size);
        assert_eq!("Size : 1.4821 MiB", 
                format!("Size : {:.4} {}", size, unit));
    }
    \end{minted}
    
    \textbf{任务3：彩色终端输出}
    
    使用 \texttt{console} crate 实现彩色输出：
    \begin{minted}{rust}
    use console::{Alignment, Term, pad_str, style};

    pub fn colored_output() {
        // INFO: 绿色
        println!("{}: Hello World!", style("INFO").green());
        
        // WARNING: 黄色、加粗、下划线
        println!(
            "{}: {}",
            style("WARNING").yellow().bold().underlined(),
            style("I'm a teapot").yellow()
        );
        
        // ERROR: 红色、加粗、居中
        let term = Term::stdout();
        let (_, width) = term.size();
        let styled_msg = style("ERROR: KERNEL PANIC!!!")
            .red().bold().to_string();
        let output = pad_str(&styled_msg, width as usize, 
                            Alignment::Center, None);
        println!("{}", output);
    }
    \end{minted}
    
    \textbf{任务4：Shape 枚举与面积计算}
    
    使用枚举实现多态形状：
    \begin{minted}{rust}
    use std::f64::consts::PI;

    enum Shape {
        Rectangle { width: f64, height: f64 },
        Circle { radius: f64 },
    }

    impl Shape {
        fn area(&self) -> f64 {
            match self {
                Shape::Rectangle { width, height } => 
                    width * height,
                Shape::Circle { radius } => 
                    radius.powi(2) * PI,
            }
        }
    }

    #[test]
    fn test_area() {
        let rectangle = Shape::Rectangle {
            width: 10.0,
            height: 20.0,
        };
        let circle = Shape::Circle { radius: 10.0 };
        
        assert_eq!(rectangle.area(), 200.0);
        assert_eq!(circle.area(), 314.1592653589793);
    }
    \end{minted}
    
    \textbf{任务5：UniqueId 唯一标识符}
    
    使用 \texttt{AtomicU16} 实现线程安全的唯一ID生成：
    \begin{minted}{rust}
    use std::sync::atomic::{AtomicU16, Ordering};

    #[derive(Debug, PartialEq)]
    struct UniqueId(u16);

    impl UniqueId {
        fn new() -> Self {
            static COUNTER: AtomicU16 = AtomicU16::new(0);
            let id = COUNTER.fetch_add(1, Ordering::SeqCst);
            UniqueId(id)
        }
    }

    #[test]
    fn test_unique_id() {
        let id1 = UniqueId::new();
        let id2 = UniqueId::new();
        assert_ne!(id1, id2);
    }
        \end{minted}

        \item \textbf{运行 UEFI Shell}
        
        使用 QEMU 启动 UEFI Shell：
        \begin{minted}[linenos=false]{bash}
        qemu-system-x86_64 -bios ./assets/OVMF.fd -net none -nographic
        \end{minted}
    
    成功进入 UEFI Shell 后，可以看到 UEFI Interactive Shell 界面，表明 OVMF 固件正确加载。

    \item \textbf{YSOS 启动！}
    
    实现了第一个 UEFI 程序（\texttt{crates/boot/src/main.rs}）：
    \begin{minted}{rust}
    #![no_std]
    #![no_main]

    #[macro_use]
    extern crate log;
    extern crate alloc;

    use core::arch::asm;
    use uefi::{Status, entry};

    #[entry]
    fn efi_main() -> Status {
        uefi::helpers::init().expect("Failed to initialize utilities");
        log::set_max_level(log::LevelFilter::Info);
        
        let std_num = "24353028";
        
        loop {
            info!("Hello World from UEFI bootloader! @ {}", std_num);
            for _ in 0..0x10000000 {
                unsafe { asm!("nop"); }
            }
        }
    }
    \end{minted}
    
    运行 \texttt{make run} 后，成功输出：
    \begin{minted}{text}
    [ INFO]: crates/boot/src/main.rs@017: 
    Hello World from UEFI bootloader! @ 24353028
    \end{minted}

\end{enumerate}

\section{思考题}

\begin{enumerate}[label=\arabic*.]
    \item \textbf{UEFI 和 Legacy（BIOS）的区别是什么？}
    
    \begin{itemize}
        \item \textbf{启动方式}：Legacy BIOS 使用 MBR 分区表，UEFI 使用 GPT 分区表
        \item \textbf{启动速度}：UEFI 支持并行初始化，启动速度更快
        \item \textbf{安全特性}：UEFI 支持 Secure Boot，可以验证引导程序签名
        \item \textbf{容量支持}：Legacy BIOS 受限于 MBR 分区表，最多支持 2TB 的硬盘
        \item \textbf{图形界面}：UEFI 支持图形化界面和鼠标操作
    \end{itemize}
    
    \item \textbf{Makefile 中的命令解释}
    
    Makefile 主要完成以下工作：
    \begin{itemize}
        \item \texttt{build}：编译 UEFI 引导程序，将生成的 \texttt{.efi} 文件复制到 \texttt{esp/EFI/BOOT/BOOTX64.EFI}
        \item \texttt{launch}：启动 QEMU，使用 OVMF 作为 UEFI 固件，挂载 esp 目录作为 FAT 文件系统
        \item \texttt{run}：先 build 再 launch
        \item \texttt{debug}：启动 QEMU 并开启 GDB 调试端口（-s -S）
    \end{itemize}
    
    \item \textbf{第三方库的源代码位置和编译时机}
    
    \begin{itemize}
        \item 源代码存储在 \texttt{\textasciitilde/.cargo/registry/src/} 目录下
        \item 编译时机：首次构建项目时，Cargo 会下载并编译所有依赖
        \item 编译结果缓存在 \texttt{target/debug/deps/} 或 \texttt{target/release/deps/}
    \end{itemize}
    
    \item \textbf{为什么使用 \texttt{\#[entry]} 而不是 main 函数？}
    
    \begin{itemize}
        \item UEFI 程序是 \texttt{no\_std} 环境，没有标准库提供的 C 运行时启动代码
        \item \texttt{\#[entry]} 宏将函数标记为程序入口点，生成必要的启动代码
        \item UEFI 规范定义了特定的入口函数签名，\texttt{\#[entry]} 宏处理这些细节
        \item 直接使用 main 函数需要操作系统提供运行时支持，而 UEFI 程序运行在裸机环境
    \end{itemize}
\end{enumerate}

\section{加分项}

\begin{enumerate}[label=\arabic*.]
    \item \textbf{彩色日志输出}
    
    基于 \texttt{log} crate 实现了自定义的彩色日志输出器，为不同日志级别设置不同颜色：
    \begin{minted}{rust}
    impl Log for LvzwLogger {
        fn log(&self, record: &Record) {
            match record.metadata().level() {
                Level::Error => println!("[{}]: {}",
                    record.level().as_str().red(), record.args()),
                Level::Warn => println!("[{}]: {}",
                    record.level().as_str().yellow(), record.args()),
                Level::Info => println!("[{}]: {}",
                    record.level().as_str().green(), record.args()),
                Level::Debug => println!("[{}]: {}",
                    record.level().as_str().blue(), record.args()),
                Level::Trace => print!("[{}], {}",
                    record.level().as_str().purple(), record.args()),
            }
        }
    }
    \end{minted}
    
    \item \textbf{简单 Shell 程序}
    
    实现了一个简易的 shell 程序，支持以下命令：
    \begin{itemize}
        \item \texttt{cd <path>}：切换当前工作目录
        \item \texttt{ls [path]}：列出目录内容
        \item \texttt{pwd}：显示当前工作目录
        \item \texttt{exit}：退出 shell
    \end{itemize}
    
    核心实现使用 \texttt{std::env::set\_current\_dir} 和 \texttt{std::fs::read\_dir}：
    \begin{minted}{rust}
    fn cd<P: AsRef<Path>>(path: P) -> ShellResult<()> {
        set_current_dir(path)?;
        pwd()?;
        Ok(())
    }

    fn ls(dir: &Path, cb: &dyn Fn(&Path)) -> ShellResult<()> {
        if dir.is_dir() {
            let mut entries = fs::read_dir(dir)?
                .map(|res| res.map(|e| e.path()))
                .collect::<Result<Vec<_>, io::Error>>()?;
            entries.sort();
            for entry in entries { cb(&entry.as_path()) }
        }
        Ok(())
    }
    \end{minted}
    
    \item \textbf{UniqueId 线程安全性验证}
    
    本实验验证了 \texttt{static mut} 和 \texttt{AtomicU16} 两种实现方式在多线程环境下的行为差异。
    
    \textbf{1. static mut 版本（不安全）}
    
    使用 \texttt{static mut} 实现计数器：
    \begin{minted}{rust}
    static mut COUNTER: u16 = 0;

    pub unsafe fn next_id_unsafe() -> u16 {
        let id = COUNTER;
        std::hint::black_box(());  // 增加竞争概率
        COUNTER += 1;
        id
    }
    \end{minted}
    
    在多线程测试中，这个版本可能产生重复 ID，因为：
    \begin{itemize}
        \item 多个线程可能同时读取 \texttt{COUNTER} 的相同值
        \item 写入操作不是原子的，可能导致数据竞争
        \item 编译器优化可能重排指令
    \end{itemize}
    
    \textbf{2. AtomicU16 版本（安全）}
    
    使用原子类型实现计数器：
    \begin{minted}{rust}
    static COUNTER: AtomicU16 = AtomicU16::new(0);

    pub fn next_id_safe() -> u16 {
        COUNTER.fetch_add(1, Ordering::SeqCst)
    }
    \end{minted}
    
    这个版本保证线程安全，因为：
    \begin{itemize}
        \item \texttt{fetch\_add} 是原子操作，读取和写入不可分割
        \item \texttt{Ordering::SeqCst} 保证顺序一致性
        \item 无需 \texttt{unsafe} 块
    \end{itemize}
    
    \textbf{3. 测试结果}
    
    \begin{minted}[linenos=false]{text}
    === 测试 static mut 版本（不安全） ===
    总 ID 数量: 4000
    唯一 ID 数量: 3847
    重复 ID 数量: 153
    检测到重复 ID！这证明了 static mut 在多线程下是不安全的。

    === 测试 AtomicU16 版本（安全） ===
    总 ID 数量: 4000
    唯一 ID 数量: 4000
    重复 ID 数量: 0
    所有 ID 都是唯一的！AtomicU16 保证了线程安全。
    \end{minted}
\end{enumerate}

\section{关于 unsafe 的思考}

Rust 的 \texttt{unsafe} 关键字是语言设计中的一个重要权衡，它既强大又危险。

\subsection*{unsafe 存在的意义}

\begin{itemize}
    \item \textbf{底层系统编程}：操作系统内核、驱动程序需要直接操作内存和硬件
    \item \textbf{FFI 互操作}：与 C/C++ 等语言的代码进行交互
    \item \textbf{基础数据结构实现}：标准库中的 \texttt{Vec}、\texttt{Box} 等需要 unsafe
    \item \textbf{性能优化}：在关键路径上绕过安全检查
\end{itemize}

\subsection*{unsafe 的风险}

\begin{itemize}
    \item 绕过 Rust 的内存安全保证
    \item 可能导致数据竞争、悬垂指针、内存泄漏
    \item 未定义行为可能在不同编译器版本间变化
    \item Bug 更难调试和复现
\end{itemize}

\subsection*{最佳实践}

\begin{enumerate}
    \item 尽量缩小 \texttt{unsafe} 块的范围
    \item 为 \texttt{unsafe} 代码提供安全的抽象接口
    \item 使用文档注释说明安全性条件
    \item 使用 Miri 等工具检测未定义行为
\end{enumerate}

\subsection*{本实验的启示}

\texttt{static mut} 在多线程下是不安全的，因为：
\begin{itemize}
    \item 多个线程可能同时读取和修改，导致数据竞争
    \item 编译器可能进行优化，导致不可预期的行为
\end{itemize}

\texttt{AtomicU16} 提供了安全的替代方案：
\begin{itemize}
    \item 使用 CPU 的原子指令保证操作的原子性
    \item 通过内存排序（\texttt{Ordering}）控制可见性
    \item 无需使用 \texttt{unsafe} 块
\end{itemize}

\textbf{结论}：\texttt{unsafe} 是 Rust 安全模型的"逃生舱"，但应该谨慎使用。在可能的情况下，优先选择安全的替代方案。Rust 的设计哲学是"安全默认，unsafe 显式"，这要求程序员在使用 \texttt{unsafe} 时必须清楚地知道自己在做什么，并承担相应的责任。

\section{实验总结}

本次实验完成了以下内容：
\begin{enumerate}
    \item 搭建了 Rust + QEMU 的开发环境
    \item 完成了 5 个 Rust 编程练习，熟悉了 Rust 的所有权、错误处理、枚举、原子操作等特性
    \item 成功运行了 UEFI Shell，理解了 UEFI 启动过程
    \item 实现了第一个 UEFI 引导程序，为后续操作系统开发奠定了基础
\end{enumerate}

\clearpage
\end{document}